\section{讨论}\label{sec:discussion}

\paragraph{为什么判定门控学不到“LLM 是否可靠”}
“LLM 在某个实例上是否站对了边”由两个量决定：LLM 表达的判断与该实例的真实状态。前者写在 LLM 的输出里，后者在闭卷条件下只能由新闻文本去推测。因此，一个读取新闻文本与 LLM 输出的门控，要预测这一目标，最直接的办法就是先从文本预测标签，再与 LLM 的判断比较。表~\ref{tab:gate-cells} 与表~\ref{tab:useful} 显示 \method{} 的判定门控与 ARG 的可用性预测器都落在这一解上。要得到不依赖标签预测的可靠性估计，需要第三个信息来源，例如外部证据，或 LLM 在同一问题上多次采样的一致性；后者是否有效留待后续工作检验。多示例学习的不对称性确实为真新闻中的声明提供了精确的监督，但这一监督同样可以由“这是一条真新闻”推出，并没有迫使模型去判断单条声明。

\paragraph{多示例假设在多大程度上成立}
式~\eqref{eq:mil} 假设假新闻中存在可以单独指认的问题声明。\ref{sec:exp-gate} 节的统计显示，在 Weibo21 上多数假新闻是由整体实例而不是某条声明触发的。这与数据的性质有关：社交媒体上的不实消息常常整体是编造的，其中每条声明单独看都“像是可能发生的”，仅凭文本无法指出哪一条有问题。反过来，真新闻如实转述不实说法的情形会违反假设的另一半；我们在提示词中要求保留归属以减少这种情况，但没有通过人工标注去度量假设被违反的比例，这是本文的一个局限。

\paragraph{关于任务本身}
风格改写与事实篡改两项测试指向同一个结论：在这两个常用数据集上，闭卷检测器学到的主要是真假新闻在写法上的差异。这与 \citet{wu2024sheepdog} 对 SLM 的观察一致，本文进一步表明它同样适用于 LLM 的闭卷判断，以及建立在这些判断之上的协同方法。把 LLM 的分析拆得更细、让它写出正反论证，改变的是分析的形式，而没有改变分析所依据的信息。对于“内容是否属实”这一问题，闭卷的 7B 模型既不掌握所需的事实，也没有表现出能从文本内部一致性中发现问题的能力。这提示：在不引入外部证据的前提下继续改进 LLM 与 SLM 的协同结构，在这类基准上能得到的主要是对文风线索更充分的利用；而声称具备“事实分析”能力的闭卷方法，应当接受风格受控的测试。

\paragraph{对实践的建议}
如果目标只是分布内精度：在算力允许时直接微调 LLM；否则在 SLM 的输出上融合一个 LLM 零样本评分，成本最低且不逊于更复杂的结构。如果系统中确实要使用 LLM 生成的分析文本，本文的结果支持为它配一个相关性检查：它几乎不损失精度，却能避免在分析与新闻不符时性能大幅下降。无论采用哪种方案，方法间 1 个点以内的差异都需要多个随机种子与显著性检验才能判断。

\paragraph{局限}
（1）分析用的 LLM 以 Qwen2.5-7B 为主，更强的模型可能掌握更多事实，结论未必外推到闭源大模型。（2）只使用了两个数据集，且都来自同一来源的时间划分。（3）风格改写与事实篡改由另一个 LLM 自动生成，事实是否保持只用数字串的保留比例做了粗略检查，没有人工核验；篡改后的句子与原文不符，但未必在现实中为假。（4）\method{} 的第一个版本曾在测试集上评估过一次，相关性门控是在看到那次结果后加入的（\ref{sec:protocol} 节）；虽然此后的选择只依据验证集，这仍然弱化了测试集的独立性。（5）本文对 ARG 的复现比其官方代码低 0.5--1 个点。（6）声明的质量、论证的质量以及定位结果都没有做人工评估。（7）\method{} 的超参数没有搜索，可能不是最优。
